\chapter{Overleaf i współpraca}\label{ch:overleaf}

Overleaf to LaTeX w przeglądarce. Nie trzeba nic instalować, projekt jest dostępny z każdego komputera, a promotor widzi tę samą wersję co Ty. Dla pracy licencjackiej pisanej pod opieką promotora to najmniej kłopotliwe środowisko. Praca lokalna też jest możliwa i opisuje ją podrozdział~\ref{sec:lokalnie}.

\section{Ścieżka studenta}\label{sec:overleafstudent}

\subsection{Konto i projekt}

Załóż konto \wyroznienie{na adres uczelniany}. Sprawdź, czy Uniwersytet udostępnia licencję instytucjonalną -- jeżeli tak, konto w domenie uczelni otrzymuje funkcje płatne. Cała ścieżka opisana niżej działa również na planie bezpłatnym.

Wzór pobierasz z repozytorium seminarium: pobierz archiwum całego repozytorium, rozpakuj je, spakuj \wyroznienie{sam katalog wzoru} do nowego archiwum i wgraj je jako nowy projekt. Alternatywnie promotor udostępni gotowy projekt wzorcowy łączem tylko do odczytu; wtedy wystarczy skopiować projekt do własnego konta.

\subsection{Ustawienia}

\begin{tabelaepi}{Ustawienia projektu na Overleaf}{tab:overleaf}
\begin{tabular}{@{}llp{6.4cm}@{}}
\toprule
Ustawienie & Wartość & Dlaczego \\
\midrule
Compiler & LuaLaTeX & kod składa się wtedy właściwym krojem pisma \\
Main document & \plik{main.tex} & inaczej kompiluje się dokument pokazowy \\
Spell check & polski & podkreśla literówki w trakcie pisania \\
\bottomrule
\end{tabular}
\end{tabelaepi}

Ustawienie kompilatora jest jednorazowe i \wyroznienie{łatwe do przeoczenia}. Objaw pominięcia: praca kompiluje się poprawnie, ale kod i nazwy funkcji mają inny krój pisma niż powinny.

\subsection{Co jest czym w projekcie}

W pliku głównym uzupełniasz metadane na górze i nic więcej. Piszesz w plikach rozdziałów. Wpisy bibliograficzne trafiają do pliku w katalogu bibliografii, a grafiki do katalogu rysunków. Dokument pokazowy służy do zaglądania, gdy nie pamiętasz zapisu.

Klasy dokumentu ani plików stylu bibliograficznego \wyroznienie{nie ruszasz}. Zmiana czegokolwiek w nich oznacza, że praca przestaje odpowiadać wymogom Instytutu, a odpowiedzialność za to spada na Ciebie, nie na wzór. Jeżeli czegoś nie da się zrobić dostępnymi poleceniami, zgłoś to promotorowi -- brakuje wtedy polecenia we wzorze.

\subsection{Kompilacja i kopie}

Pierwsza kompilacja po dodaniu nowego powołania albo etykiety może pokazać znaki zapytania; po drugiej znikną. Rozwijane menu obok przycisku kompilacji zawiera polecenie pełnego przebiegu od zera -- użyj go, gdy dokument przestaje odpowiadać zmianom albo gdy bibliografia nie chce się odświeżyć.

Plan bezpłatny ma limit czasu kompilacji. Praca licencjacka mieści się w nim spokojnie, ale gdy zaczniesz się o niego ocierać, zakomentuj w pliku głównym włączenia rozdziałów, nad którymi akurat nie pracujesz. Przed oddaniem odkomentuj wszystkie.

Overleaf przechowuje projekt na swoich serwerach, ale \wyroznienie{to nie jest kopia zapasowa}. Raz w tygodniu pobierz źródła projektu i zapisz je poza komputerem, na którym pracujesz; przed każdym większym przemeblowaniem pracy dodatkowo. Historia zmian na planie bezpłatnym jest ograniczona i kopii zapasowej nie zastępuje.

\subsection{Udostępnienie i oddanie}

Udostępnij projekt promotorowi na jego adres, z prawem edycji. Na planie bezpłatnym każdy projekt można udostępnić jednej osobie -- dokładnie tyle, ile trzeba. Recenzentowi, gdy zajdzie potrzeba, udostępnij łącze tylko do odczytu; łącza nie liczą się do limitu współpracowników.

Do Archiwum Prac trafia plik wynikowy pobrany z projektu. Przed pobraniem wykonaj pełną kompilację od zera, sprawdź brak nierozwiązanych powołań i odsyłaczy oraz przejdź listę kontrolną z podrozdziału~\ref{sec:listakontrolna}.

\section{Plik bibliograficzny a Zotero}

Dodatek Better BibTeX odświeża plik bibliograficzny \wyroznienie{na Twoim dysku}. Overleaf o tym nie wie, więc plik trzeba do projektu przenieść. Najprostszy sposób to wgranie ręczne: w drzewie plików wskaż katalog bibliografii, prześlij zaktualizowany plik i potwierdź nadpisanie. Wystarczająco dobre, jeżeli robisz to raz na tydzień. Overleaf potrafi też wczytać publiczną bibliotekę Zotero i odświeżać ją na żądanie, co jest wygodne, ale wymaga ustawienia biblioteki jako publicznej. Przy dostępie do synchronizacji z repozytorium plik odświeża się razem z resztą projektu.

Niezależnie od sposobu: \wyroznienie{po każdej podmianie pliku uruchom pełną kompilację od zera}. Bibliografia budowana jest z pliku pomocniczego, który potrafi zostać na starej wersji.

\section{Ścieżka promotora}

\subsection{Rozdanie wzoru grupie}

Najmniej kłopotliwy układ, działający również na kontach bezpłatnych, wygląda tak. Prowadzący zakłada \wyroznienie{jeden projekt wzorcowy} z zawartości katalogu wzoru i włącza w nim łącze tylko do odczytu. Rozdaje to jedno łącze wszystkim uczestnikom. Każda osoba kopiuje projekt do własnego konta, otrzymując niezależną kopię, i udostępnia ją prowadzącemu z prawem edycji.

Zalety układu: łącze do odczytu nie zużywa limitu współpracowników, poprawka we wzorcu nie psuje niczyjej pracy, a prowadzący ma dostęp do wszystkich projektów z jednego pulpitu. Wadą jest to, że poprawka we wzorcu nie trafia do kopii już wykonanych -- dlatego \wyroznienie{wzór zamyka się przed pierwszymi zajęciami}, a późniejsze zmiany rozsyła jako opis, co podmienić.

\subsection{Przegląd prac}

Bez funkcji płatnych do dyspozycji pozostaje zwykła edycja i komentarze w treści. Przy kilkunastu pracach warto ustalić jedną konwencję i trzymać się jej przez cały rok. Konwencja proponowana w seminarium to uwagi zapisane jako komentarze LaTeX-a, które \wyroznienie{nie pojawiają się w gotowym dokumencie}:

\begin{zapis}
%% UWAGA: teza z tego akapitu nie wynika z przywolanego zrodla.
%% UWAGA: brakuje powolania.
%% PYTANIE: skad wartosc progu 0,7?
\end{zapis}

Uwaga stoi wtedy dokładnie przy problematycznym zdaniu, student widzi ją przy pisaniu, a gotowy dokument pozostaje czysty. Wyszukanie wszystkich uwag w projekcie to jedno kliknięcie w polu wyszukiwania.

Zasada zamykająca: \wyroznienie{student usuwa komentarz dopiero po naniesieniu poprawki}. Pusty wynik wyszukiwania oznacza, że wszystkie uwagi zostały zaadresowane; ten punkt znajduje się na liście kontrolnej przed oddaniem.

Przy licencji instytucjonalnej dochodzą śledzenie zmian, komentarze w panelu bocznym i pełna historia wersji -- wtedy powyższa konwencja staje się zbędna.

\subsection{Kontrola formalna}

Kontrolę formalną prowadź \wyroznienie{na złożonym dokumencie, nie w źródle}. Marginesy, interlinia, numeracja stron i układ elementów końcowych widać dopiero w składzie. Zestawienie wymogów wraz ze wskazaniem, co realizuje klasa, a co pozostaje po stronie autora, zawiera rozdział~\ref{ch:wymogi}.

W źródle warto sprawdzić dwie rzeczy: czy student nie zmieniał klasy dokumentu ani plików stylu bibliograficznego, co załatwia porównanie z wzorcem, oraz czy w projekcie nie zostały komentarze robocze i uwagi.

Przy kilkunastu pracach przegląd na żądanie nie działa. Sprawdzony układ to stałe okna: oddanie fragmentu do ustalonej daty, przegląd w ciągu tygodnia, omówienie na zajęciach. Fragmenty oddawane w kolejności pisania, a nie w kolejności rozdziałów -- kolejność podaje tabela~\ref{tab:kolejnosc}.

\section{Praca lokalna w Positronie}\label{sec:lokalnie}

Skoro pakiet R budujesz w Positronie, możesz w nim też pisać pracę. Kod i tekst są wtedy w jednym środowisku, a wykresy trafiają prosto do katalogu rysunków. Potrzebujesz dystrybucji TeX-a oraz rozszerzenia do LaTeX-a.

\begin{lstlisting}[language=bash,numbers=none]
cd wzor-pracy
latexmk -lualatex main.tex
\end{lstlisting}

Zaletami są brak limitu czasu kompilacji, praca bez dostępu do sieci i pełna historia w repozytorium. Wadami -- konieczność zainstalowania i utrzymania dystrybucji TeX-a oraz to, że udostępnienie promotorowi wymaga repozytorium zamiast jednego kliknięcia. Układ mieszany bywa najwygodniejszy: pakiet i wykresy lokalnie, tekst na Overleaf.

Repozytorium pakietu i pracę trzymaj \wyroznienie{osobno}. Praca zawiera pliki, które nie mają nic wspólnego z pakietem, a repozytorium pakietu ma być czyste -- to ono jest oceniane jako projekt dyplomowy i to jego adres podajesz na stronie tytułowej. Jeżeli prowadzisz pracę w repozytorium, ignoruj artefakty kompilacji.

\section{Problemy specyficzne dla Overleaf}

\begin{tabelaepi}{Problemy w pracy na Overleaf}{tab:problemyoverleaf}
\begin{tabular}{@{}p{5cm}p{4.5cm}p{4.5cm}@{}}
\toprule
Objaw & Przyczyna & Naprawa \\
\midrule
kod bez właściwego kroju pisma & kompilator ustawiony na pdfLaTeX & zmień kompilator na LuaLaTeX \\
kompiluje się dokument pokazowy & zły dokument główny & wskaż plik główny pracy \\
bibliografia nie odświeża się & stary plik pomocniczy & pełna kompilacja od zera \\
kompilacja przerwana limitem czasu & za dużo naraz przy szkicowaniu & zakomentuj nieużywane włączenia \\
grafika nie znaleziona & plik poza katalogiem rysunków & przenieś plik, podawaj samą nazwę \\
polskie znaki jako znaki zastępcze & plik w innym kodowaniu & zapisz ponownie w UTF-8 \\
zmiany współpracownika niewidoczne & stara karta przeglądarki & odśwież stronę \\
projekt nie chce się skopiować & łącze bez prawa kopiowania & poproś o łącze do odczytu \\
\bottomrule
\end{tabular}
\end{tabelaepi}

Katalog błędów kompilacji niezależnych od Overleaf zawiera podrozdział~\ref{sec:diagnostyka}.
